\documentclass{beamer}

% Packages
\usepackage{amsmath}
\usepackage{amsfonts}
\usepackage{amssymb}
\usepackage{graphicx}
\usepackage{subfig}
\usepackage[utf8]{inputenc}
\usepackage{booktabs}
\usepackage{bm}
\usepackage{xcolor}

% -----------------------------------------------------------------------------
% THEME & COLOR CUSTOMIZATION
% -----------------------------------------------------------------------------
\usetheme{CambridgeUS}

% Define Custom Purple
\definecolor{DeepPurple}{RGB}{80, 0, 80} 

% Overwrite CambridgeUS Red colors with Purple
\setbeamercolor{palette primary}{bg=DeepPurple,fg=white}
\setbeamercolor{palette secondary}{bg=DeepPurple!80!black,fg=white}
\setbeamercolor{palette tertiary}{bg=DeepPurple!60!black,fg=white}
\setbeamercolor{palette quaternary}{bg=DeepPurple!40!black,fg=white}

% Structure (Bullets, Table of Contents, etc.)
\setbeamercolor{structure}{fg=DeepPurple}
\setbeamercolor{section in toc}{fg=black} 
\setbeamercolor{subsection in toc}{fg=DeepPurple}

% Titles
\setbeamercolor{title}{bg=DeepPurple,fg=white}
\setbeamercolor{frametitle}{fg=DeepPurple, bg=white}

% Blocks (Theorems, Definitions)
\setbeamercolor{block title}{bg=DeepPurple, fg=white}
\setbeamercolor{block body}{bg=white, fg=black}
\setbeamercolor{block title alerted}{bg=black, fg=white}
\setbeamercolor{block body alerted}{bg=white, fg=black}
\setbeamercolor{block title example}{bg=white, fg=DeepPurple}

% Itemize/Enumerate bullets
\setbeamercolor{item projected}{bg=DeepPurple, fg=white}
\setbeamercolor{itemize item}{fg=DeepPurple}
\setbeamercolor{itemize subitem}{fg=DeepPurple}

% Remove navigation symbols for a cleaner look
\setbeamertemplate{navigation symbols}{}

% Add frame numbers to footline
\setbeamertemplate{footline}[frame number]{}

% -----------------------------------------------------------------------------
% METADATA
% -----------------------------------------------------------------------------
\title[Krylov Matrix Exponentials]{On Krylov Subspace Approximations to the Matrix Exponential Operator}
\subtitle{Reproducing Results from Hochbruck \& Lubich (1997)}
\author{Purple Team \\[0.5em] Marek Fürst, Ján Kovačovský, Ida Nyárfás, Jaroslav Paidar, Petr Rašek}
\date{\today}

\begin{document}

% -----------------------------------------------------------------------------
% TITLE SLIDE
% -----------------------------------------------------------------------------
\begin{frame}
    \titlepage
\end{frame}

% -----------------------------------------------------------------------------
% TABLE OF CONTENTS
% -----------------------------------------------------------------------------
\begin{frame}{Outline}
    \tableofcontents
\end{frame}

% -----------------------------------------------------------------------------
% INTRODUCTION
% -----------------------------------------------------------------------------
\section{Introduction}

\begin{frame}{Motivation}
    \begin{itemize}
        \item \textbf{Goal}: Compute $e^{\tau A}v$ efficiently for large sparse matrices $A$.
        \item \textbf{Context}: Stiff ODEs; exponential integrators vs. linear solves $(I - \gamma \tau A)^{-1}v$.
        \item \textbf{Key Insight}: Krylov for $e^{\tau A}v$ converges faster; depends on $\mathcal{F}(A)$.
    \end{itemize}
\end{frame}

% TODO: add plots with spectrum (eigenvalue distribution)
% TODO: add comparison between Arnoldi and Lanczos case

% -----------------------------------------------------------------------------
% THEORY: ARNOLDI
% -----------------------------------------------------------------------------
\section{Arnoldi Approximation}

\begin{frame}{Arnoldi Process and Approximation}
    \begin{block}{Krylov Subspace}
        $\mathcal{K}_m(A, v) = \text{span}\{v, Av, \dots, A^{m-1} v\}$.
    \end{block}
    
    \begin{block}{Relation}
        $A V_m = V_m H_m + h_{m+1,m} v_{m+1} e_m^T$, $H_m = V_m^* A V_m$ (Hessenberg).
    \end{block}

    \begin{block}{Approximation}
        $e^{\tau A}v \approx \beta V_m e^{\tau H_m} e_1$, $\beta = \|v\|_2$.
    \end{block}
\end{frame}

\begin{frame}{Error Bound (Lemma 1)}
    For convex $E \supset \mathcal{F}(A)$, error:
    \[
        \| e^{\tau A}v - \beta V_m e^{\tau H_m} e_1 \| \le \frac{M}{2\pi} \int_{\Gamma} |e^{\tau \lambda} - q_{m-1}(\lambda)| \cdot |\phi(\lambda)|^{-m} \, |d\lambda|.
    \]
    \begin{itemize}
        \item $\phi$: Conformal map; decay vs. $e^{\tau \lambda}$ growth.
    \end{itemize}
\end{frame}

% -----------------------------------------------------------------------------
% RESULTS: SYMMETRIC NEGATIVE SEMIDEFINITE
% -----------------------------------------------------------------------------
\section{Symmetric Negative Semidefinite}

\begin{frame}{Case 1 (Theorem 2)}
    Eigenvalues in $[-4\rho, 0]$.
    \begin{alertblock}{Bounds}
        \begin{itemize}
            \item $\sqrt{4\rho\tau} \le m \le 2\rho\tau$: $\epsilon_m \le 10 e^{-m^2 / (5\rho\tau)}$.
            \item $m \ge 2\rho\tau$: $\epsilon_m \le 10 (\rho\tau)^{-1} e^{-\rho\tau} (e\rho\tau / m)^m$.
        \end{itemize}
    \end{alertblock}
    Superlinear at $m \approx \sqrt{\rho \tau}$.
\end{frame}

\begin{frame}{Results}
    \begin{figure}
        \centering
        \includegraphics[width=0.85\textwidth]{../graphics/negative_semidefinite.pdf}
        \caption{Error vs. CG for $(I - A)^{-1}v$; eigenvalues $[-40, 0]$, $N=1001$.}
    \end{figure}
\end{frame}

% -----------------------------------------------------------------------------
% RESULTS: SKEW-HERMITIAN
% -----------------------------------------------------------------------------
\section{Skew-Hermitian}

\begin{frame}{Case 2 (Theorem 4)}
    Eigenvalues on imaginary axis, length $4\rho$.
    \begin{block}{Bound}
        $m \ge 2\rho\tau$: $\epsilon_m \le 12 e^{-(\rho\tau)^2/m} (e\rho\tau / m)^m$.
    \end{block}
    Delay until $m \approx \rho\tau$.
\end{frame}

\begin{frame}{Results}
    \begin{figure}
        \centering
        \includegraphics[width=0.85\textwidth]{../graphics/skew_hermitian.pdf}
        \caption{Error vs. BiCG for $(I - A)^{-1}v$; eigenvalues $[-20i, 20i]$, $N=1001$.}
    \end{figure}
\end{frame}

% -----------------------------------------------------------------------------
% LANCZOS METHODS
% -----------------------------------------------------------------------------
\section{Lanczos Method}

\begin{frame}{Lanczos Approximation}
    \begin{itemize}
        \item Biorthogonal $V_m, W_m$; $H_m = D_m^{-1} W_m^* A V_m$ (tridiagonal for Hermitian).
        \item Matches Arnoldi for symmetric/skew cases.
        \item Bounds involve pseudospectra.
    \end{itemize}
\end{frame}

\begin{frame}{Results}
    \begin{columns}
        \begin{column}{0.5\textwidth}
            \begin{figure}
                \includegraphics[width=\textwidth]{../graphics/lanczos_negative_semidefinite.pdf}
                \caption{Symmetric negative semidefinite.}
            \end{figure}
        \end{column}
        \begin{column}{0.5\textwidth}
            \begin{figure}
                \includegraphics[width=\textwidth]{Lanczos_skew_Hermitian.png}
                \caption{Skew-Hermitian.}
            \end{figure}
        \end{column}
    \end{columns}
\end{frame}

% -----------------------------------------------------------------------------
% CONCLUSION
% -----------------------------------------------------------------------------
\section{Conclusion}

\begin{frame}{Conclusion}
    \begin{itemize}
        \item Reproduced bounds/plots for symmetric negative semidefinite and skew-Hermitian.
        \item Hierarchy: Faster for symmetric ($m \approx \sqrt{\rho\tau}$) than skew ($m \approx \rho\tau$).
        \item Exponential Krylov superior for stiff ODEs.
    \end{itemize}
\end{frame}

\begin{frame}{Reference}
    \begin{itemize}
        \item M. Hochbruck and C. Lubich, \emph{On Krylov Subspace Approximations to the Matrix Exponential Operator}, SIAM J. Numer. Anal., 34(5):1911--1925, 1997.
    \end{itemize}
\end{frame}

\end{document}
